\documentclass[11pt,a4paper]{article}
\usepackage[utf8]{inputenc}
\usepackage[T1]{fontenc}
\usepackage[italian]{babel}
\usepackage{amsmath,amssymb,graphicx,booktabs,siunitx,hyperref}
\usepackage{listings}
\lstdefinestyle{faust}{basicstyle=\ttfamily\small,breaklines=true,
  keepspaces=true,columns=fullflexible,frame=single,framesep=4pt}
\lstset{style=faust}
\usepackage[margin=2.5cm]{geometry}

\title{ltglide: il glissando del tone-burst di Linkwitz\\\large Diario di studio --- porting C\texttt{++} e plugin per la suite SEAM-LTM}
\author{Giuseppe Silvi --- SEAM}
\date{2026}

\begin{document}
\maketitle
\begin{abstract}
Questo documento \`e il diario di studio di \emph{ltglide}, il generatore di glissando del tone-burst sagomato di Linkwitz.
Racconta il porting in C\texttt{++} della specifica Faust \texttt{seam.linkwitz.lib} (prefisso \texttt{slw}), gi\`a discussa in dettaglio nel diario di \emph{ltburst}, e la costruzione del plugin che la ospita.
La teoria del burst sagomato e la derivazione della specifica Faust restano nel diario di \emph{ltburst} (\texttt{plugins/ltburst/doc/study/ltburst-study.tex}), che questo documento richiama per riferimento.
\end{abstract}
\tableofcontents

\section{Il glissando come generalizzazione del burst fisso}\label{sec:generalizzazione}

Il tone-burst sagomato di Linkwitz~\cite{linkwitz1980shaped} fissa il numero di cicli $N=5$ e lascia libera la frequenza portante $f_0$, ottenendo un segnale a Q costante che \emph{ltburst} riproduce fedelmente a frequenza singola.
\emph{ltglide} generalizza questo schema facendo scorrere la frequenza portante lungo una passata, cos\`i lo strumento diventa un rilievo continuo esteso su tutta la banda coperta dallo sweep.
Uno sweep esterno, calcolato dalla funzione \texttt{sweepfreq} della libreria \texttt{seam.linkwitz.lib}, produce il progresso di frequenza in funzione del tempo di passata.
Un motore a grana consuma quello sweep con un campionamento-e-tenuta a ogni proprio inizio, cos\`i ciascuna grana resta a frequenza singola per tutta la propria durata.
Il risultato \`e una sequenza di burst di Linkwitz via via pi\`u acuti o pi\`u gravi, ciascuno internamente identico a un burst di \emph{ltburst} alla propria frequenza tenuta.
Questa costruzione rende \emph{ltglide} uno strumento di misura continuo: la stessa banda-costante di un singolo burst si ripete lungo tutto lo spettro coperto dalla passata.

La funzione \texttt{sweepfreq} offre due leggi di percorrenza, gi\`a specificate in Faust e riportate fedelmente nel core C\texttt{++} (\texttt{ltglide\_dsp.h}):
\begin{equation}
\texttt{sweepfreq}(f_0,f_1,\mathrm{smode},p) =
\begin{cases}
f_0 + (f_1-f_0)\,p, & \mathrm{smode}=0 \ (\text{lineare}) \\
f_0 \left(\dfrac{f_1}{f_0}\right)^{p}, & \mathrm{smode}=1 \ (\text{esponenziale})
\end{cases}
\end{equation}
dove $p \in [0,1]$ \`e il progresso della passata.
La legge lineare percorre hertz uguali per unit\`a di progresso, e il suo punto medio $p=0.5$ vale la media aritmetica $(f_0+f_1)/2$.
La legge esponenziale percorre ottave uguali per unit\`a di progresso, e il suo punto medio vale la media geometrica $\sqrt{f_0 f_1}$.
La suite sceglie l'esponenziale come modalit\`a predefinita, coerente con la percezione logaritmica dell'altezza e con la scala delle frequenze gi\`a usata dai controlli \texttt{F0}/\texttt{F1} del plugin.

\section{Il motore a grana retriggerata (Approccio A)}\label{sec:motore}

Il cuore di \emph{ltglide} \`e \texttt{GlissBurst}, il porting diretto della funzione Faust \texttt{glissburst} gi\`a derivata e discussa nel diario di \emph{ltburst}~(\S5.3, Approccio A).
Il motore porta avanti due stati accoppiati campione per campione: una rampa di grana \texttt{phase\_} in $[0,1)$ e una frequenza tenuta \texttt{fg\_}.
Ogni campione avanza la rampa di un incremento pari al reciproco del periodo di grana corrente, calcolato dalla frequenza tenuta del campione precedente.
Quando la rampa supera l'unit\`a, oppure al primissimo campione della passata, il motore dichiara un nuovo inizio di grana.
Proprio a ogni inizio di grana la frequenza tenuta \texttt{fg\_} si aggiorna al valore corrente dello sweep esterno \texttt{fsig}, e resta fissa a quel valore per tutta la durata della grana successiva: questo \`e il campionamento-e-tenuta descritto nella Sezione~\ref{sec:generalizzazione}.
La portante e la finestra di Hann derivano entrambe dalla stessa coppia \texttt{(phase\_, fg\_)} appena aggiornata, cos\`i condividono un unico riferimento di fase e restano sincrone per costruzione.

\begin{lstlisting}
glissburst(N,delta,dmode,fsig) = sin(2*ma.PI*u) * win
  with {
    phase = grain : _,!;
    fg    = grain : !,_;
    grain = loop ~ (_,_)
      with {
        loop(pphase,pfg) = nphase, nfg
          with {
            den    = max(20.0, pfg);
            Tg     = select2(dmode, max(delta, N/den), N/den + delta);
            inc    = 1.0/max(1.0, Tg*ma.SR);
            adv    = pphase + inc;
            onset  = (adv >= 1.0) | (1 - 1');
            nphase = adv - floor(adv);
            nfg    = select2(onset, pfg, fsig);
          };
      };
    den = max(20.0, fg);
    Tg  = select2(dmode, max(delta, N/den), N/den + delta);
    u   = fg * phase * Tg;
    win = (u < N) * (0.5 - 0.5*cos(2*ma.PI*u/N));
  };
\end{lstlisting}

Il porting C\texttt{++} riproduce questa struttura ricorsiva in modo identico: \texttt{GlissBurst::process(fsig)} calcola prima la coppia successiva \texttt{(phase\_, fg\_)} a partire dai valori precedenti, esattamente come il blocco \texttt{loop} Faust, e solo dopo usa quella coppia aggiornata per sintetizzare portante e finestra nello stadio di uscita.
Una conseguenza diretta di questa struttura riguarda il primissimo campione della passata: la rampa avanza di un incremento \emph{prima} di essere usata nello stadio di uscita dello stesso campione, cos\`i il primo campione della grana resta a una piccola distanza misurabile dallo zero-crossing.
Il test \texttt{doctest} verifica infatti che il primo campione resti entro $3.8 \times 10^{-6}$ dallo zero, con un confronto a tolleranza assoluta, perch\'e una tolleranza relativa attorno a un valore atteso esattamente nullo degenera a una tolleranza pressoch\'e nulla.
Questa piccola distanza dallo zero \`e un'eredit\`a fedele della definizione ricorsiva Faust, e differisce dalla convenzione a contatore esatto (\texttt{c} come stato) usata da \emph{ltburst}, gi\`a discussa e motivata nel suo diario di studio.

Una conseguenza numerica di questa costruzione lega \emph{ltglide} a \emph{ltburst} in modo verificabile.
A frequenza costante il motore a grana retriggerata degenera a un generatore a periodo fisso, e il modo gap con $f=\SI{1000}{\hertz}$, $\delta=\SI{0.3}{\second}$ produce inizi di grana distanziati di $N/f+\delta = 0.305\,\mathrm{s}$, cio\`e $14640$ campioni a $\SI{48}{\kilo\hertz}$.
Questo \`e esattamente il periodo del burst fisso di \emph{ltburst} a $f_0=\SI{1000}{\hertz}$ con \emph{dwell} $=\SI{0.3}{\second}$, gi\`a validato nel suo diario e nel suo record di validazione.
Il numero coincide per costruzione, perch\'e entrambi i motori condividono la stessa formula di periodo quando la frequenza \`e costante, e questa coincidenza numerica costituisce un'ancora di verifica tra i due oggetti pur restando due implementazioni C\texttt{++} indipendenti, ciascuna con codice proprio nel rispettivo plugin.

\section{I due modi di temporizzazione e la guardia di sicurezza}\label{sec:timing}

Il motore offre due leggi per il periodo di grana $T_g$, gi\`a specificate in Faust e riprodotte identiche nel porting C\texttt{++} tramite il parametro \texttt{Timing} del plugin.
Il modo \emph{passo} fissa la distanza inizio-inizio a $T_g = \max(\delta, N/f_g)$, cos\`i il passo tra le grane resta costante e il silenzio residuo si accorcia quando la frequenza tenuta sale.
Il modo \emph{gap} fissa il silenzio fine-inizio a $T_g = N/f_g + \delta$, cos\`i la distanza inizio-inizio cresce quando la frequenza tenuta scende, perch\'e la grana stessa dura pi\`u a lungo.
Entrambe le leggi condividono la stessa guardia di sicurezza $\mathrm{den} = \max(20\,\mathrm{Hz}, f_g)$, che protegge il calcolo di $T_g$ quando la frequenza tenuta \`e ancora nulla, al primissimo inizio della passata, oppure quando lo sweep raggiunge il proprio estremo grave.
Questa guardia rende il modo passo sicuro contro la sovrapposizione delle grane anche al limite inferiore dello sweep, perch\'e a $\SI{20}{\hertz}$ una grana di cinque cicli dura gi\`a $\SI{250}{\milli\second}$, un valore che il termine $\max(\delta, N/f_g)$ rispetta sempre per costruzione.

\section{Il transport della passata}\label{sec:transport}

Il plugin colloca il motore a grana dentro una passata delimitata da due impulsi Dirac, gestita dalla classe \texttt{GlideTransport} e indipendente dal motore stesso.
Un trigger avvia la passata con un impulso di testa, seguito da un tratto di silenzio (\emph{lead}) di cinque secondi, che lascia assestare l'ambiente di misura prima che lo sweep abbia inizio.
Segue la fase di glissando vera e propria, della durata impostata dal parametro \texttt{Sweep Time}, durante la quale il progresso $p$ cresce da zero a uno e alimenta \texttt{sweepfreq} e quindi il motore a grana.
Al termine dello sweep un secondo tratto di silenzio (\emph{tail}) di cinque secondi precede l'impulso di coda, simmetrico all'impulso di testa e utile a delimitare la registrazione della passata.
Il parametro \texttt{Loop}, se attivo, fa seguire all'impulso di coda un'attesa di due secondi e poi un nuovo impulso di testa, cos\`i la passata si ripete in modo automatico e deterministico.
Un nuovo trigger \`e ignorato mentre una passata \`e gi\`a in corso, cos\`i il transport mantiene una passata alla volta e la sua struttura resta prevedibile per uno strumento di misura.

Il loop serve principalmente alla fase di ricevitore, ancora da progettare (fase 3b-ii della suite), dove pi\`u passate consecutive verranno mediate per migliorare il rapporto segnale-rumore della misura.
Gli impulsi Dirac di testa e di coda offrono a quel ricevitore un riferimento temporale netto per allineare e segmentare le passate prima della media, un ruolo analogo a quello che gli stessi impulsi svolgono gi\`a come marcatori di inizio e fine registrazione per un operatore umano.

\section{Rinvii}\label{sec:rinvii}

La misura degli spettri della passata, il confronto con le figure originali di Linkwitz e l'analisi delle risposte in frequenza ricavate dal glissando restano compiti della fase di ricevitore.
Questo diario si ferma alla validazione del motore a grana e del transport come generatore standalone, gi\`a documentata nel record di validazione (\texttt{plugins/ltglide/doc/ltglide-validation.md}).
La teoria del tone-burst sagomato, la derivazione della specifica Faust e la discussione della differenza tra le convenzioni di stato di \emph{ltburst} e di \emph{ltglide} restano nel diario di \emph{ltburst}, che questo documento richiama come fonte primaria.

\bibliographystyle{plain}
\bibliography{refs}
\end{document}
